\documentclass[12pt]{article}
\input{iati_structure.tex}
\renewcommand{\bottomtitlespace}{3cm}

\usepackage[english]{babel}

\begin{document}

\renewcommand{\headerLang}{Russian}
\renewcommand{\problemName}{AB13. Vision}

\renewcommand{\headerLeft}{IATI Day 1, Senior group}
\renewcommand{\headerRightFirst}{XVII INTERNATIONAL ADVANCED TOURNAMENT IN INFORMATICS}
\renewcommand{\headerRightSecond}{ONLINE 2026}

\renewcommand{\tl}{$5$ s}
\renewcommand{\ml}{$1024$ MB}
\problem{Task \problemName}
\addauthor{Author: Emil Indzhev}{2026}{04}{16}

Космический корабль USS Enterprise (с Сашкой в качестве нового капитана) путешествует по Вселенной, представляющей собой бесконечную двумерную сетку секторов. В каждом секторе $(X, Y)$\footnote{$X$ -- номер строки (увеличивающийся вниз), а $Y$ — номер столбца (увеличивающийся вправо).} находится маяк с силой видения $V_{X, Y}$. Находясь в этом секторе, корабль использует маяк для сканирования всех секторов $(X', Y')$ с $X - V_{X, Y} \leq X' \leq X + V_{X, Y}$ и $Y - V_{X, Y} \leq Y' \leq Y + V_{X, Y}$. Таким образом, видимые секторы образуют квадрат со стороной длиной $2V_{X, Y} + 1$. Для каждого такого сектора единственная доступная информация -- это сила видения маяка в нём, т.е. $V_{X', Y'}$ (это связано с тем, что маяки выступают одновременно в роли передатчиков и приемников). После сканирования корабль телепортируется в один из этих видимых секторов, поскольку маяки также выполняют функцию телепортаторов.

К сожалению, у маяков есть один существенный недостаток: они могут перевернуть ось $X$ сканирования, или ось $Y$, или обе оси (то есть существует четыре возможных ориентации сканирования). Обратите внимание, что корабль использует это возможно перевернутое сканирование, чтобы решить, в какой сектор телепортироваться, но при этом телепортация также использует ту же ориентацию, например, если ось $Y$ перевёрнута и корабль телепортируется влево (уменьшая ось $Y$) сканирования, он фактически телепортируется вправо. Это означает, что корабль действительно телепортируется в сектор, выбранный при сканировании.

Кроме того, корабль не обладает памятью, поэтому его решение о том, куда телепортироваться, должно зависеть только от результатов сканирования маяка. Также стратегия корабля должна быть детерминированной — получив одинаковые результаты сканирования, он должен выбрать для телепортации тот же сектор, что и на сканировании.

Корабль начинает движение из неизвестных координат и всегда (независимо от начальной точки и ориентации получаемых сканирований) должен иметь возможность двигаться в направлении произвольно увеличивающихся $X$ и $Y$ (можно сказать, что в конечном итоге он должен достичь сектора $(X, Y)$ такого, что $X, Y \geq 10^{100}$). Однако для этого решающее значение имеют как стратегия перемещения корабля, так и расположение сил видения во Вселенной.

Вот тут-то вы и вступаете в дело -- вам нужно разработать прямоугольный шаблон $N$ на $M$ с силами видения, который будет бесконечно повторяться во Вселенной: $V_{X, Y} = P_{X \bmod N, Y \bmod M}$. Вам также нужно разработать стратегию перемещения корабля, то есть функцию, которая принимает данные сканирования, полученные от маяка, и выбирает, в какой из этих секторов телепортироваться дальше.

Поскольку маяки с более высокой силой видения стоят дорого, ваша цель -- получить корректное решение, стараясь при этом минимизировать среднюю силу видения в вашем шаблоне, то есть среднее значение $P$.

Поскольку эта задача может показаться довольно сложной, вы решаете начать с тренировочной задачи: той же задачи, но в одномерном пространстве. В этой более простой задаче мы отбрасываем измерение $X$ и используем только измерение $Y$. Таким образом, ваш шаблон представляет собой просто последовательность длиной $M$, которая повторяется следующим образом: $V_Y = P_{Y \bmod M}$. Сканирования, создаваемые маяками, также представляют собой последовательности длиной $2V_Y + 1$ (содержащие сектора $Y'$ такие, что $Y - V_{X, Y} \leq Y' \leq Y + V_{X, Y}$). Здесь существует только две возможных ориентации сканирования (обычная или перевёрнутая), и корабль должен быть способен достичь $Y \geq 10^{100}$.

Ваши решения для одномерного и двумерного случаев будут оцениваться отдельно, и вы можете получить баллы за один случай, даже если у вас нет правильного решения для другого.

\subsection{Детали реализации}

Вам надо реализовать 4 функции: \verb|getVisionPattern1d|, \verb|getMove1d|, \verb|getVisionPattern2d| и \verb|getMove2d|. Первые две используются для одномерного случая, а вторые две — для двумерного. В любом тесте будет вызвана только одна из этих пар функций, но для корректной работы решения (для компиляции) необходимо реализовать все четыре (даже если реализации одной пары пусты).

\begin{lstlisting}
 std::vector<int> getVisionPattern1d()
\end{lstlisting}

Эта функция будет вызвана один раз, если тест проводится для одномерного случая. Она должна возвращать последовательность $P$ длиной $M$ -- шаблон силы одномерного видения, который необходимо повторить. $M$ и все элементы $P$ должны находиться в диапазоне от $1$ до $60$.

\begin{lstlisting}
 int getMove1d(std::vector<int> v)
\end{lstlisting}

Эта функция будет вызываться один раз за каждое возможное сканирование, если проверка проводится для одномерного случая. Ей передается сканирование, полученное от маяка, в виде последовательности длиной $2V + 1$ (где $V$ — интенсивность видения маяка в центральном секторе последовательности). Она должна вернуть индекс $T$ в последовательности, в который должен телепортироваться корабль, так чтобы $0 \leq T < 2V + 1$.

\begin{lstlisting}
 std::vector<std::vector<int>> getVisionPattern2d()
\end{lstlisting}

Эта функция будет вызвана один раз, если проверка проводится для двумерного случая. Она должна возвращать прямоугольную таблицу $P$ размером $N$ на $M$ — шаблон интенсивности двумерного видения, который бедт должен повторяться. Значения $N$, $M$ и всех элементов таблицы $P$ должны находиться в диапазоне от $1$ до $60$.

\begin{lstlisting}
 std::pair<int, int> getMove2d(std::vector<std::vector<int>> v)
\end{lstlisting}

Эта функция будет вызываться один раз за каждое возможное сканирование, если проверка проводится в двумерном случае. Ей передается сканирование, полученное от маяка, в виде квадратной таблицы размером $2V + 1$ на $2V + 1$ (где $V$ — сила видения маяка в центральном секторе квадрата). Функция должна возвращать, в какой сектор должен телепортироваться корабль, в виде пары индексов $S$ и $T$ в таблице, таких что $0 \leq S, T < 2V + 1$.

Для заданной размерности $D$ ($1$ или $2$) теста грейдер должен убедиться, что ваша программа генерирует корректный шаблон -- что она делает корректный выбор телепортации для всех возможных сканирований и что ваши решения работают независимо от начальных координат и последовательности ориентаций сканирований.

\subsection{Ограничения}

\vspace{0.1em}

\begin{itemize}
\item $1 \leq D \leq 2$
\item $1 \leq N, M, V_{X, Y}, V_Y \leq 60$
\end{itemize}

\subsection{Подзадачи и система оценки}

\begin{table}[H]
\begin{tblr}{|Q[c,m]|Q[c,m]|Q[c,m]|Q[c,m]|}
\hline
\textbf{Подзадача} & \textbf{Баллы} & \textbf{$D$} & \textbf{$A^*$} \\
\hline
$1$ & $24$ & $1$ & $1.5$ \\
\hline
$2$ & $76$ & $2$ & $1.125$ \\
\hline
\end{tblr}
\end{table}
\FloatBarrier

Ваш балл за данную подзадачу (если ваше решение верное) зависит от средней силы видения в вашем шаблоне $P$ для этой подзадачи. Пусть это значение будет $A$, а целевым значением для подзадачи — $A^*$. Ваш балл (доля баллов за подзадачу) будет следующим:

$$1 - {\left(1 - \frac{A^* - 1}{\max{\left(A, A^*\right)} - 1}\right)}^{0.8}$$

\subsection{Пример шаблона}

Предположим, ваше решение возвращает следующий шаблон с $N=3$ и $M=2$:

\[
\left[
\begin{array}{c c}
1 & 2 \\
3 & 4 \\
5 & 6
\end{array}
\right]
\]

Теперь рассмотрим возможные варианты сканирования в секторе $(0, 1)$, сила видения в котором составляет $2$. Существует $4$ возможных ориентации (некоторые из которых дают одинаковый результат сканирования на этом шаблоне):

\[
\begin{array}{cc}
\text{Ориентация $0$: Нет переворотов} & \text{Ориентация $1$: Ось $Y$ перевёрнута} \\
\left[
\begin{array}{ccccc}
4 & 3 & 4 & 3 & 4 \\
6 & 5 & 6 & 5 & 6 \\
2 & 1 & 2 & 1 & 2 \\
4 & 3 & 4 & 3 & 4 \\
6 & 5 & 6 & 5 & 6
\end{array}
\right]
&
\left[
\begin{array}{ccccc}
4 & 3 & 4 & 3 & 4 \\
6 & 5 & 6 & 5 & 6 \\
2 & 1 & 2 & 1 & 2 \\
4 & 3 & 4 & 3 & 4 \\
6 & 5 & 6 & 5 & 6
\end{array}
\right]
\\
\\
\text{Ориентация $2$: Ось $X$ перевёрнута} & \text{Ориентация $3$: Обе оси перевёрнуты} \\
\left[
\begin{array}{ccccc}
6 & 5 & 6 & 5 & 6 \\
4 & 3 & 4 & 3 & 4 \\
2 & 1 & 2 & 1 & 2 \\
6 & 5 & 6 & 5 & 6 \\
4 & 3 & 4 & 3 & 4
\end{array}
\right]
&
\left[
\begin{array}{ccccc}
6 & 5 & 6 & 5 & 6 \\
4 & 3 & 4 & 3 & 4 \\
2 & 1 & 2 & 1 & 2 \\
6 & 5 & 6 & 5 & 6 \\
4 & 3 & 4 & 3 & 4
\end{array}
\right]
\end{array}
\] 

Поскольку $0$ и $1$ идентичны, и поскольку корабль детерминирован, для них будет сделан только один вызов функции \verb|getMove2d|. Предположим, ваше решение возвращает ход $(0, 1)$ при этом сканировании. При ориентации $0$ это будет означать перемещение на $1$ влево и на $2$ вверх, а при ориентации $1$ — перемещение на $1$ вправо и на $2$ вверх.

\subsection{Пример грейдера}

Первый входной параметр грейдера -- $D$. После этого он выполнит все вызовы вашей функции: сначала получит шаблон $P$, а затем получит все варианты телепортации. Затем он начнёт взаимодействие с консолью, запросив начальные координаты корабля. После этого он выведет текущее состояние: координаты и «сырое» (без переворотов) сканирование. Затем он запросит используемую ориентацию ($0$ — без переворотов, $1$ — переворот оси $Y$, $2$ — переворот оси $X$ и $3$ — переворот обеих осей). После этого он выведет сканирование в заданной ориентации, а также выбранный кораблем вариант движения. Корабль будет перемещён, и процесс повторится. Обратите внимание, что пример грейдера не будет автоматически проверять правильность вашего решения.


\end{document}
